\chapter{Limiti e sviluppi futuri}
\label{chap:discussione}

Questo capitolo accoppia ogni limite misurato allo sviluppo che lo
affronterebbe provando a risolverlo. Quando uno sviluppo tocca la sicurezza, indica anche che cosa
servirebbe perché le garanzie del Paragrafo\ref{sec:ch3:sicurezza} continuino a
valere.

\section{La leggibilità dei dati grezzi ed il modello di riconoscimento}
\label{sec:ch6:ml-riconoscimento}

Il riconoscimento della sorgente, descritto nel
Paragrafo\ref{subsec:ch3:riconoscimento}, funziona per convenzioni note, quindi è
verificabile ma non copre nulla che non sia già stato registrato. Un file con
etichette che non si possono interpretare senza la convenzione del laboratorio che
le ha prodotte viene rifiutato inoltre un formato di ingresso nuovo richiederebbe un
intervento sul codice. Questo comportamento è corretto perché indovinare
riprodurrebbe l'errore del Paragrafo\ref{sec:ch3:angoli} senza un modello clinico
a fianco con cui accorgersene, però limita i dati che il sistema può leggere.

Lo sviluppo proposto è un modello di machine learning che deduca l'identità
anatomica di un punto dalla geometria e dal movimento, indipendentemente dal nome
assegnato da chi ha prodotto il file. L'informazione necessaria è già nei dati, le
distanze fra punti dello stesso segmento sono stabili nel tempo, il corpo è
approssimativamente simmetrico ed un punto distale si muove in modo diverso da uno prossimale. Anche il materiale
di addestramento esiste già, cioè le sessioni convertite con la corrispondenza
registrata fra etichetta originale e punto canonico.

Le conseguenze di questo sviluppo sarebbero due. Se l'identità dei punti non dipende dai nomi,
aggiungere un formato significa solamente saperne leggere i numeri integrando la
libreria adatta, quindi i formati diffusi nell'animazione, le uscite dei sistemi
markerless ed altre librerie di stima di posa diventerebbero sorgenti ugualmente utilizzabili.
Inoltre poter leggere qualunque archivio pubblico di pose cambierebbe di molto la scala dei
dati disponibili.

Il vincolo da mantenere è che un riconoscimento appreso può sbagliare 
quindi il modello va affiancato al riconoscimento per convenzioni, non lo
sostituisce. Il modello propone un'assegnazione con la propria confidenza, sopra
una soglia alta la proposta viene sottoposta a chi carica, mentre sotto la soglia
il file resta rifiutato. I punti identificati in questo modo porterebbero inoltre
una provenienza propria.

\section{La località del software ed il lancio in rete}
\label{sec:ch6:mcp-online}

Oggi server e database funzionano in locale ed una sessione si carica indicandone
il nome in una cartella ammessa perché il server legge solo il proprio
filesystem. Lo sviluppo naturale è un server raggiungibile in rete che accetti la
sessione grezza inserita direttamente nella chat, dove l'utente allega il file.
L'agente lo passa allo strumento, la pipeline lo converte e lo carica, la sessione
diventa interrogabile nella stessa conversazione. Questo lavoro riguarda ogni
garanzia esistente, oltre al trasporto.

\paragraph{Ricezione del contenuto.} Il file arriverebbe come flusso di byte, che
andrebbe confinato in un'area temporanea con limiti su dimensione, frequenza e
formato. Poiché un file malformato può far terminare la libreria nativa che legge
il C3D, la conversione resterebbe inoltre isolata in un sottoprocesso.

\paragraph{Identità di chi carica.} Servirebbero l'autenticazione degli agenti ed
una nozione di proprietà delle sessioni. Con chiavi uniche nell'archivio, due
utenti che caricano file con lo stesso nome si sovrascriverebbero, quindi una
chiave di sessione opaca generata al caricamento, con il nome originale
conservato come metadato non visibile all'agente, diventa un prerequisito. 
Risolve anche il limite di pseudonimizzazione della chiave attuale.

\paragraph{Il database in rete.} Le garanzie locali restano valide però non
bastano più, perché cadono tre assunzioni. La prima è che chi raggiunge il
database sia legittimo, in locale accedere alla porta significa accedere alla
macchina mentre in rete la password del ruolo diventa l'unica barriera. La seconda
è che i dati siano pochi e di prova, invece un archivio condiviso contiene persone
reali ed il legame fra pseudonimo ed identità diventa un bersaglio. La terza è che
il volume delle richieste sia benigno, mentre il costo della ricerca completa
presuppone un solo interrogante alla volta. La Tabella\ref{tab:ch6:sicurezza}
associa a ciascuna minaccia la sua contromisura.

\begin{table}[htbp]
\centering
\caption{Contromisure per un'esposizione in rete associate alla minaccia che
affrontano.}
\label{tab:ch6:sicurezza}
\begin{tabularx}{\textwidth}{XXl}
\toprule
\textbf{Minaccia} & \textbf{Contromisura} \\
\midrule
Accesso non autorizzato alla porta del database &
  Il database non è esposto: solo il server applicativo lo raggiunge, in rete
  privata \\
Credenziale applicativa compromessa &
  Autenticazione a certificato client, rotazione periodica, segreti da un
  gestore esterno \\
Agente non identificato &
  Autenticazione degli agenti e proprietà delle sessioni \\
Esaurimento di risorse &
  Limiti sulla frequenza delle richieste e sul costo per interrogazione \\
Correlazione fra pseudonimo e identità &
  Chiave di sessione opaca generata al caricamento \\
Compromissione dell'host &
  Cifratura a riposo del volume & da fare \\
Assenza di tracciabilità &
  Registro degli accessi per agente e per sessione \\
\bottomrule
\end{tabularx}
\end{table}

La scrittura dalla conversazione, che in locale tiene la credenziale di scrittura
nel processo del server, in rete ne diventerebbe la parte più esposta. La strada
corretta è separare i due processi, con un servizio di caricamento che possiede
la propria credenziale e che il server di interrogazione può soltanto sollecitare.

\paragraph{Un'interfaccia web.} Con il software online diventa utile anche
un'interfaccia web per caricare le sessioni e gestire i soggetti senza passare
dalla chat, compresa la cancellazione su richiesta dei dati di una persona. Anche
questa interfaccia richiederebbe autenticazione ed un accesso all'anagrafica
riservato ad un ruolo da amministratore.

\section{Gli sport e le condizioni coperti dagli alert}
\label{sec:ch6:sport}

Gli intervalli di letteratura coprono oggi due sport, cammino e corsa, e
descrivono la condizione in piano. Il sistema conosce il gesto ma non la
condizione, quindi applica gli stessi intervalli anche alle prove in pendenza, che
sono quattro delle sei condizioni del dataset di laboratorio\cite{vielemeyer2026}.

Su una prova in piano di un soggetto, ad esempio, \texttt{list\_alerts} non
riporta nessuno scostamento, mentre sulla prova dello stesso soggetto su rampa al
10\% ne riporta sette, tutti su serie misurate con copertura 1.00, come mostra la
Tabella\ref{tab:ch6:rampa}.

\begin{table}[htbp]
\centering
\caption{Gli scostamenti su una singola prova su rampa al 10\%. La stessa prova
in piano dello stesso soggetto non ne produce nessuno. Valori in gradi.}
\label{tab:ch6:rampa}
\begin{tabular}{llrrl}
\toprule
\textbf{Giunto} & \textbf{Metrica} & \textbf{Osservato} &
  \textbf{Atteso} & \textbf{Direzione} \\
\midrule
ginocchio sx & escursione  & 78.60 & 50--75  & sopra \\
ginocchio sx & picco max   & 78.56 & 50--75  & sopra \\
ginocchio dx & escursione  & 76.75 & 50--75  & sopra \\
anca sx      & escursione  & 28.88 & 30--60  & sotto \\
anca dx      & escursione  & 29.28 & 30--60  & sotto \\
anca sx      & picco min   & $-4.60$ & $-25$--$-5$ & sopra \\
caviglia sx  & picco max   & 20.86 & 5--20   & sopra \\
\bottomrule
\end{tabular}
\end{table}

Gli scostamenti non vanno tutti nella stessa direzione, il ginocchio flette più
dell'atteso e l'anca meno, che è esattamente il profilo di un cammino in pendenza.
Sull'intero archivio l'andamento è lo stesso, con 1.64 segnalazioni per sessione
in piano, 4.62 su rampa al 7.5\% e 6.21 al 10\%, come già visto nel
Paragrafo\ref{subsec:ch4:demo-statistica}. Sulle rampe quindi gli avvisi
descrivono il protocollo più che il soggetto, perché il sistema applica un
riferimento ad un caso che quel riferimento non copre, e lo dichiara nella fonte
di ogni avviso.

Per questo limite gli sviluppi possibili sono due. Il primo è estendere il file dei riferimenti con
un ambito per condizione, come la pendenza, e con nuovi gesti accompagnati dai
relativi dati di letteratura come è già stato fatto per cammino e corsa. Nella
struttura questo è semplice, però richiede valori di letteratura per ogni
condizione, che è un lavoro di dominio.

Il secondo è affiancare ai valori fissi un secondo livello di sessioni di
riferimento, come le
sessioni precedenti della stessa persona, oppure un insieme di sessioni per uno
sport non coperto dalla letteratura. L'infrastruttura esiste già, ma ci sono delle condizioni da rispettare. 
Servono un ambito ed un'etichetta propri, così che non si confonda con un
riferimento di letteratura;  serve un numero minimo di
sessioni, perché una linea di base fatta di una sola prova coincide con la prova stessa.


\section{Il visualizzatore esterno e la visualizzazione in chat}
\label{sec:ch6:visualizzazione}

Oggi il visualizzatore apre una finestra nativa sulla macchina che ospita il
server, fuori dalla conversazione, e questo ha senso quando chi interroga e chi
ospita sono la stessa persona ma non per un agente remoto. Lo sviluppo è una
visualizzazione dentro la chat, restituita come immagine o come vista
interattiva, con la possibilità di andare direttamente ai fotogrammi voluti, ad
esempio quelli in cui compare un alert. Il visualizzatore sa già scrivere un
fotogramma su file senza aprire finestre e sa già calcolare in quali fotogrammi si
produce uno scostamento, quindi lo sviluppo consiste nel portare queste capacità direttamente 
nella risposta dello strumento.

\section{Difetti noti degli strumenti e correzioni previste}
\label{sec:ch6:difetti}

\paragraph{Il ciclo interrogato non è scelto per rappresentatività.} La ricerca di
cicli simili usa il ciclo indicato, che per impostazione è il primo, e su una
sessione con molta variabilità questa scelta conta. La correzione è proporre il
ciclo più vicino alla mediana della sessione, oppure avvisare quando la
dispersione fra i cicli è alta.

\paragraph{Gli alert usano l'intera prova.} Le metriche giudicate al caricamento
sono calcolate sull'intera registrazione, compresi i fotogrammi fuori dai cicli.
Su una prova di laboratorio con due soli cicli, ad esempio, il ginocchio sinistro
arriva a $-12.3$ gradi sulla prova ed a $-10.7$ sul ciclo, ed uno degli
scostamenti cade prima dell'inizio del primo ciclo. La correzione coincide con la
valutazione per ciclo indicata sopra, con la sessione riportata come sintesi.

\paragraph{Escursioni oltre 180 gradi non segnalate.} Nelle riprese da telefono
alcuni angoli derivati superano i $\pm180$ gradi e producono escursioni
fisicamente impossibili. Nella sessione \texttt{running\_\allowbreak 20260220\_\allowbreak 121922}, ad
esempio, il tronco arriva a 185 gradi e la spalla sinistra a 212, e né la
conversione né gli strumenti lo rilevano. La correzione è svolgere gli angoli prima del
calcolo delle metriche ed aggiungere un controllo di plausibilità che segnali con
un avviso le escursioni fuori dai limiti fisici.

\paragraph{La ripetibilità mescola i cicli dei due lati.} Quando i cicli vengono
dal ripiego sulla spalla, la sessione contiene cicli segmentati sull'una e
sull'altra spalla, spesso sovrapposti, e lo strumento
\texttt{session\_repeatability} li confronta tutti per ogni giunto, quindi una parte della distanza misura in realtà
lo sfasamento fra i due lati. La correzione è confrontare per ogni giunto solo i
cicli del lato corrispondente, o almeno dichiarare nella risposta da quale lato vengono.

\paragraph{Il gesto può essere scelto dall'agente.} Il caricamento rifiuta una
sorgente senza gesto, però il controllo riguarda la chiamata, quindi un agente che
deduce il gesto e lo passa allo strumento carica senza chiedere conferma, come è
successo con entrambi gli agenti nel caso 1 del
Paragrafo\ref{subsec:ch4:caso1}. Il server non può sapere chi abbia scelto il
gesto, però può renderlo visibile riportandolo nella risposta come gesto da
confermare quando la sorgente non lo dichiara, così che l'agente debba
presentarlo all'utente.

\section{I limiti dei dati}
\label{sec:ch6:dati}

\paragraph{Le riprese da telefono non contengono l'arto inferiore.} Nelle riprese
usate in questo lavoro l'inquadratura esclude le gambe, quindi quelle sessioni
sono dati di tronco ed arto superiore. Il sistema dichiara le assenze, però non può
recuperare ciò che la ripresa non contiene. Lo sviluppo è un protocollo di ripresa
con il soggetto intero nel fotogramma, ad una distanza tale che occupi circa un
terzo dell'immagine con i piedi dentro il bordo.

\paragraph{Il modello di segmento.} Con i soli centri articolari la caviglia porta
uno scostamento sistematico di circa dodici gradi in assenza di una prova statica,
quindi le escursioni restano confrontabili mentre i valori assoluti no. L'anca
conserva un guadagno di scala di circa 1.25, perché il bacino nasce dai centri
articolari e non dalle spine iliache, ed il varo-valgo del ginocchio non è
osservabile. Tutti questi limiti sono dichiarati nel dato, e lo sviluppo è usare
marker aggiuntivi o una prova statica quando la sorgente li fornisce.

